\bigskip
\noindent Autres pistes de l'onglet Findings pour ce thème, dans des cotes qui n'ont pas de section ici :\par

\begin{encadre}
\begin{theoreme*}[Grothendieck, pages 11 à 16, polarisations hodgiennes\piste{10-hodgian-polarisations-torsor}]\nonverifie
Soient $G$ un groupe algébrique réel muni d'homomorphismes $\mathbb G_m \xrightarrow{i} G \xrightarrow{\varepsilon} \mathbb G_m$, $\eta = i(-1)$ et $G' = \mathrm{Ker}\,\varepsilon$. Deux éléments $C, C'$ de $G'(\mathbb R)$ de carré $\eta$ qui définissent des polarisations de $\mathrm{Rep}(G)$ sont, après conjugaison par $G'(\mathbb R)^\circ$, liés par $C' = zC$ avec $z$ central d'ordre divisant~$2$, et $z \neq 1$ donne une autre polarisation : les polarisations hodgiennes forment un torseur sous un groupe d'éléments centraux de $2$-torsion. Sur toute $\otimes$-catégorie graduée sur $K \subset \mathbb R$, les $\otimes$-automorphismes $z$ de l'identité tels que $z^2 = 1$, triviaux sur $K(1)$, opèrent librement sur les polarisations.
\end{theoreme*}
\end{encadre}
{\small Les feuillets en esquissent une démonstration incomplète (pages 15 et 16). La conjugaison de la page~15 tient en trois lignes abrégées où quatre mots sont illisibles, et le groupe qui opère, la $2$-torsion de $Z \cap G'$, est le choix de l'édition. Mathlib n'a pas les catégories tannakiennes ni leurs polarisations. On se reportera aux notes de Grothendieck (\lot{10}{1}{11}).\par}

\begin{encadre}
\begin{theoreme*}[Grothendieck, pages 62 à 65, caractère numérique et tore de rang un\piste{10-numerical-character-rank-one-torus}]\nonverifie
Une catégorie tannakienne sur un corps $k$ dont le lien est un tore de dimension~$1$ est déterminée à équivalence près par son caractère numérique : l'ensemble des classes d'objets simples, les constantes de structure de $K(\mathcal M)$, les centres $Z_\sigma$ des algèbres d'endomorphismes avec leurs classes de Brauer, et le rang. Pour le tore $G$ des éléments de norme~$1$ d'une extension quadratique $Z/k$, l'application $H^2(k, G) \to H^2(k, R_{Z/k}\mathbb G_m) = \mathrm{Br}(Z)$ est en effet injective, par Hilbert~90.
\end{theoreme*}
\end{encadre}
{\small La page passe par \q{on vérifie aussitôt} l'équivalence entre le caractère numérique et les données $(G, u_\sigma(\xi))$. Mathlib n'a pas les catégories tannakiennes ni leurs gerbes. On se reportera aux notes de Grothendieck (\lot{10}{4}{62}).\par}

\begin{encadre}
\begin{theoreme*}[Grothendieck, pages 10 à 12, groupe de Picard de l'espace de Grothendieck--Springer\piste{123-grothendieck-springer-pic-twisted-torus}]\nonverifie
Soient $S$ un schéma normal intègre, $G$ un $S$-schéma en groupes réductif dont le tore maximal abstrait $T$ n'est pas déployé, et $\Gamma$ l'image de $\pi_1(S)$ dans $\mathrm{Aut}\,X^*(T)$. L'application $\mathrm{Pic}(\mathrm{Bor}) \to \mathrm{Pic}(\tilde G)$ est injective, de conoyau contenu dans $H^1(\Gamma, X^*(T))$, et bijective quand ce $H^1$ est nul. Un fibré en droites sur $\mathrm{Bor}$ est ample si et seulement si son image réciproque sur $\tilde G$ l'est, sans hypothèse de normalité sur~$S$.
\end{theoreme*}
\end{encadre}
{\small La condition $H^1(\Gamma, X^*(T)) = 0$ vient d'une marge et de la page~12, lues avec plusieurs mots incertains ou illisibles ; la réciproque pour $S$ régulier, que la marge annonce, n'est pas reprise ici. Mathlib n'a pas les schémas en groupes réductifs ni leurs schémas de Borel. On se reportera aux notes de Grothendieck (\lot{123}{1}{10}).\par}

\begin{encadre}
\begin{theoreme*}[Grothendieck, pages 21, 28, 30 et 34, rigidifications de Kostant\piste{123-kostant-rigidifications-count}]\nonverifie
Soient $G$ un groupe simple adjoint de nombre de Coxeter $h$ premier à la caractéristique, et $z$ l'ordre du groupe $\mathfrak z$ tel que $N(T) \cap N(T')$ soit extension de $(\mathbf Z/h)^*$ par $\mathbf Z/h \times \mathfrak z$. Le groupe $G$ opère librement sur les quadruplets $(T, B, T', B')$ formés de tores maximaux $T, T'$ en apposition au sens de Kostant et de sous-groupes de Borel $B \supset T$, $B' \supset T'$. Le quotient est fini étale de rang $\mathrm{card}(W)^2/(h\,z\,\varphi(h))$, et de rang $\mathrm{card}(W)^2/(h\,z)$ quand on se donne aussi un générateur de $T \cap N(T') \simeq \mu_h$.
\end{theoreme*}
\end{encadre}
{\small La lettre à Kostant du 22~octobre~1969 (page~21) et la page~34 donnent ce compte sans démonstration. L'édition le déduit de la théorie de Kostant et de Springer \cite{kostant1959,springer1974}, par un argument que Lean ne vérifie pas. Soit $x = \rho^\vee(e^{2i\pi/h})$ l'élément principal de $T$ ; un $t \in T$ qui normalise $T'$ envoie l'élément cyclique $E'$ sur un multiple de lui-même, d'où $\alpha_i(t) = \theta(t)^{-1} = \lambda$ avec $\lambda^h = 1$, et $T \cap N(T') = \langle x \rangle$. L'image de $N(T) \cap N(T')$ dans $W$ est formée des $w$ tels que $w(x) \in \langle x \rangle$ : le stabilisateur de $x$ est le groupe $\Omega \simeq P^\vee/Q^\vee$, d'ordre $z$, car le centralisateur de $x$ dans le groupe adjoint n'est pas connexe, et $x$ n'est conjugué à $x^k$ que pour $k$ premier à $h$ ; chacun de ces $w$ se relève, par la conjugaison des sous-espaces de Cartan des $\mathfrak g^{\zeta^j}$ sous $Z_G(x)^\circ = T$. D'où l'ordre $h\,z\,\varphi(h)$, et le compte, puisque le stabilisateur d'un quadruplet est $T \cap T' = 1$. Un comptage par ordinateur sur $\mathbf F_q$ donne 2, 4 et 12 orbites pour $A_2$, $B_2$ et $G_2$. Reste la question que la lettre pose à Kostant : $B$ et $B'$ sont-ils toujours en position générale ? C'est vrai pour $A_2$, où les vecteurs propres de $E'$ sont les vecteurs de Fourier. Mathlib n'a pas les tores maximaux ni les sous-groupes de Borel des groupes algébriques. On se reportera aux notes de Grothendieck (\lot{123}{2}{21}).\par}

\begin{encadre}
\begin{theoreme*}[Grothendieck, pages 10 à 13, pas de résolution de Steinberg sans changement de base\piste{123-steinberg-no-resolution-picard}]\nonverifie
Sur une base $S$, il n'existe pas de résolution du morphisme de Steinberg $G \to T/W$ qui soit un isomorphisme au-dessus de $G^{\mathrm{reg}}$ et dont le lieu exceptionnel soit de codimension $\geq 2$ fibre par fibre : son groupe de Picard serait fini modulo $\mathrm{Pic}(S)$, alors que celui de l'espace de Grothendieck--Springer $\tilde G$ contient le réseau des poids et son cône ample. Le changement de base $T \to T/W$ est ainsi rendu nécessaire par une comparaison de groupes de Picard, et non par la monodromie.
\end{theoreme*}
\end{encadre}
{\small Les feuillets en esquissent une démonstration incomplète (pages 10 à 13). La phrase de conclusion s'interrompt, et l'argument demande une résolution régulière et projective sur $G$, hypothèses que la page ne formule pas. Mathlib n'a pas les schémas en groupes réductifs ni leurs groupes de Picard. On se reportera aux notes de Grothendieck (\lot{123}{1}{10}).\par}

\begin{encadre}
\begin{theoreme*}[Grothendieck, pages 71 à 77, catégorie de Picard de Brandt\piste{15-brandt-picard-category}]\nonverifie
Soit $C$ une catégorie de Deuring sur un anneau de Dedekind, équivalente à celle des modules inversibles sur un ordre $\mathfrak o$ d'une algèbre $E$ dont tous les automorphismes sont intérieurs et dont les bimodules inversibles contenus dans $E$ commutent entre eux. Le $2$-groupe $\mathrm{Eqv}(C)$ des autoéquivalences porte une structure symétrique strictement commutative, indépendante de l'objet choisi pour l'identifier aux bimodules. Les catégories de Deuring d'un même type sont donc des catégories-torseurs sous une même catégorie de Picard stricte, la catégorie de Picard de Brandt.
\end{theoreme*}
\end{encadre}
{\small C'est l'édition qui affirme, d'après Auslander--Goldman et Reiner, que la condition de commutation vaut pour les ordres maximaux. Mathlib n'a pas les catégories de Picard ni les catégories-torseurs. On se reportera aux notes de Grothendieck (\lot{15}{4}{71}).\par}

\begin{encadre}
\begin{theoreme*}[Grothendieck, pages 116 à 118, foncteurs fibres réels\piste{15-real-fibre-functor-criterion}]\nonverifie
Soit une catégorie motivique sur $\mathbf Q$ de lien commutatif $G$, de caractère de Tate $\varepsilon$ et de cocaractère des poids $j$ avec $\varepsilon j(\lambda) = \lambda^2$, et soit $U = \mathrm{Ker}(\varepsilon|_{G^\circ})$, où $U^\circ$ est un tore compact. La catégorie admet un foncteur fibre à valeurs dans les espaces vectoriels réels si et seulement si $U$ est connexe ; sinon $U \simeq U^\circ \times \mu_2$. Sa sous-catégorie de poids pairs, de groupe $G/j(\mu_2)$, en admet donc toujours un.
\end{theoreme*}
\end{encadre}
{\small Les feuillets en esquissent une démonstration incomplète (page 116). La condition intermédiaire~(ii) des deux propositions est en partie illisible, l'édition la rétablit, et l'implication qui y passe est la sienne. Mathlib n'a pas les catégories tannakiennes ni les foncteurs fibres. On se reportera aux notes de Grothendieck (\lot{15}{6}{116}).\par}

\begin{encadre}
\begin{theoreme*}[Grothendieck, pages 27 à 36, semi-continuité des rangs\piste{27-semicontinuity-proof}]\nonverifie
Pour un schéma en groupes lisse, séparé et de présentation finie sur une base quelconque, la dimension de la partie affine et les rangs nilpotent, affine-nilpotent et unipotent sont semi-continus supérieurement ; le rang abélien, le rang abélien plus réductif et le rang semi-simple sont semi-continus inférieurement.
\end{theoreme*}
\end{encadre}
{\small SGA~3, X~8.7, énonce ces inégalités sans démonstration. La réduction à un trait suppose les invariants localement constructibles, propriété que ni la page ni la lecture ne vérifient, et l'équivalence entre les inégalités (3) et (9) que la page affirme n'est pas justifiée. Mathlib n'a pas les schémas en groupes lisses ni leur décomposition de Chevalley. On se reportera aux notes de Grothendieck (\lot{27}{2}{27}).\par}

\begin{encadre}
\begin{theoreme*}[Grothendieck, pages 10 à 14, liens sur un corps de Tsen\piste{45-bands-tsen-field}]\nonverifie
Sur un corps $C_1$, tout lien algébrique, lisse ou non, provient d'un $k$-schéma en groupes, et $H^2(k, L)$ se réduit alors à sa classe neutre.
\end{theoreme*}
\end{encadre}
{\small Les feuillets en esquissent une démonstration incomplète (pages 11 à 14), par deux principes de réduction pour les gerbes de réalisations. Le deuxième cas suppose une conjugaison des sous-groupes de Cartan infinitésimaux qui n'est pas établie, et le troisième n'est pas lu. Mathlib n'a pas les liens ni la cohomologie non abélienne de degré~$2$. On se reportera aux notes de Grothendieck (\lot{45}{1}{10}).\par}

\begin{encadre}
\begin{theoreme*}[D'après les pages 3 à 10, forme corrigée par l'édition, surjectivité en $H^1$ en dimension un\piste{45-h1-surjectivity-dim-one}]\nonverifie
Soit $k$ un corps tel que $\mathrm{Br}(k') = 0$ pour toute extension finie séparable $k'/k$. Pour tout épimorphisme $G \to G'$ de schémas en groupes algébriques sur $k$, non nécessairement lisses ni linéaires, à noyau non nécessairement commutatif, l'application $H^1(k, G) \to H^1(k, G')$ de cohomologie plate est surjective.
\end{theoreme*}
\end{encadre}
{\small La page suppose $\mathrm{cd}(k) \leq 1$ ; l'hypothèse sur les groupes de Brauer, la dimension $\leq 1$ au sens de Serre, est de l'édition. Les feuillets en esquissent une démonstration incomplète (pages 3 à 9). Le lemme~4 invoque une conjugaison des sous-algèbres de Cartan qui n'est pas établie et qui tombe en défaut en général en caractéristique~$p$, et l'énoncé de la page~9 conclut que $H^1(k,G) \to H^1(k,G')$ est un épimorphisme pour tout épimorphisme $G \to G'$ de groupes algébriques sur~$k$. Mathlib n'a pas la cohomologie plate des schémas en groupes. On se reportera aux notes de Grothendieck (\lot{45}{1}{3}).\par}

\begin{encadre}
\begin{theoreme*}[Grothendieck, page 13, principe d'identité et lemme~B\piste{55-lemma-b-identity-principle}]\nonverifie
Soit $X$ un espace annelé quasi-compact, de faisceau structural cohérent et d'anneaux locaux noethériens, qui vérifie un principe d'identité : pour tout Idéal cohérent $J$ et tout point $y$ de $Y = \mathrm{Supp}(\mathcal O/J)$ où $\mathcal O_{Y,y}$ est intègre, il existe un voisinage ouvert $U$ de $y$ dans $Y$ sur lequel toute section nulle près de $y$ est nulle. Alors toute suite croissante de sous-faisceaux cohérents d'un faisceau cohérent est stationnaire, et $\Gamma(X, \mathcal O_X)$ est noethérien.
\end{theoreme*}
\end{encadre}
{\small La page énonce le lemme~B et son corollaire sans démonstration, et propose les polycylindres analytiques compacts comme cas d'application ; le passage au corollaire par le lemme~A est de la lecture. Mathlib n'a pas les faisceaux cohérents sur un espace annelé, et le pas central, celui où sert le principe d'identité, n'est pas vérifié en Lean ; l'édition ne vérifie pas non plus que les conditions de la page suffisent sur un espace annelé quelconque. Lean vérifie ce qui entoure ce pas, énoncé ci-dessous. On se reportera aux notes de Grothendieck (\lot{55}{1}{13}).\par}

\begin{encadre}
\begin{theoreme*}[D'après la page 13, forme réduite par l'édition, le lemme~B moins le principe d'identité\piste{55-lemma-b-identity-principle}]\verifie{Folder55}
Soient $X$ un espace quasi-compact et $(S_k)_{k \in \mathbf N}$ une suite croissante de sous-faisceaux d'un préfaisceau d'ensembles sur $X$. On suppose qu'en chaque point $x$ les germes des $S_k$ se stabilisent, et que cette stabilisation se propage uniformément : pour tout $x$ et tout $k$, il existe un voisinage ouvert $W$ de $x$ tel que tout $S_l$, $l \geq k$, qui a en $x$ les mêmes germes que $S_k$, coïncide avec $S_k$ sur $W$. Alors la suite $(S_k)$ est stationnaire. Sans l'hypothèse d'uniformité, la conclusion tombe en défaut. Enfin, si les idéaux de type fini $I$ d'un anneau $A$ s'envoient de façon croissante sur des objets $\widetilde I$ d'une classe où toute suite croissante est stationnaire, avec $\Gamma(\widetilde I) = I$, alors $A$ est noethérien.
\end{theoreme*}
\end{encadre}
{\small \emph{Démonstration.} Un sous-faisceau est ici une famille d'ensembles de sections, stable par restriction et locale : une section y appartient dès qu'elle y appartient localement. En chaque point $x$, on prend l'indice $k_x$ où les germes se stabilisent, puis le voisinage $W_x$ que donne l'uniformité pour $k_x$ : sur $W_x$, tous les $S_l$, $l \geq k_x$, coïncident avec $S_{k_x}$. Par quasi-compacité, un nombre fini de $W_x$ recouvrent $X$ ; soit $K$ le plus grand des $k_x$ correspondants. Pour $l \geq K$, une section de $S_l$ sur un ouvert $U$ appartient, sur chaque $U \cap W_x$, à $S_l = S_{k_x} \subset S_K$, donc à $S_K$ par localité : la suite est stationnaire à partir de $K$. La stabilisation des germes est ce que donne l'hypothèse~(iv) : la fibre d'un faisceau cohérent est un module de type fini sur l'anneau noethérien $\mathcal O_{X,x}$, où toute suite croissante de sous-modules est stationnaire. Pour le contre-exemple, on prend le compactifié d'Alexandrov $\mathbf N \cup \{\infty\}$, le faisceau de toutes les fonctions à valeurs dans $\{0, 1\}$, et pour $S_k$ les fonctions qui ne valent $1$ qu'aux points $0, \ldots, k-1$. En un entier $n$, les germes sont constants à partir de $k = n + 1$ ; en $\infty$, une section de $S_k$ est nulle sur le voisinage ouvert obtenu en retirant les points $0, \ldots, k-1$, et les germes sont tous ceux de $S_0$. Pourtant la fonction indicatrice de $\{k\}$ est dans $S_{k+1}$ et pas dans $S_k$ : les sous-faisceaux s'accumulent en $\infty$, ce que le principe d'identité est là pour interdire. Pour le corollaire, une suite croissante d'idéaux de type fini $I_k$ donne la suite croissante des $\widetilde{I_k}$, stationnaire, et $I_k = \Gamma(\widetilde{I_k})$ l'est aussi ; la condition des chaînes ascendantes sur les idéaux de type fini suffit à rendre $A$ noethérien. Les noms sont \texttt{stationary\_of\_locally\_stationary} (le pas de quasi-compacité), \texttt{stationary\_of\_germs} (le lemme~B sous l'hypothèse d'uniformité), \texttt{germs\_stabilize\_of\_noetherian} (le rôle de (iv)), \texttt{not\_stationary\_without\_uniformity} (le contre-exemple) et \texttt{isNoetherianRing\_of\_tilde} (le corollaire), dans \href{\leanblob/Folder55.lean}{\texttt{Folder55.\allowbreak{}lean}}. Ne sont pas démontrés : que le principe d'identité~(v) entraîne l'hypothèse d'uniformité, ce qui demanderait les faisceaux cohérents, les supports des $\mathcal O/\mathcal J$ et une décomposition primaire dans les fibres, et que le lemme~A fournisse $I \mapsto \widetilde I$ avec $\Gamma(\widetilde I) = I$. Le principe~(v) ne parle que des sections sur un voisinage $U$ fixé, pas sur les ouverts plus petits où vivent les sous-faisceaux qui s'accumulent ; l'édition ne sait pas s'il suffit sur un espace annelé quelconque.\par}
